\chapter{Sleep Apnea}

\section*{Definition}
Sleep apnea is a sleep-related breathing disorder characterized by recurrent episodes of partial (hypopnea) or complete (apnea) upper airway blockage during sleep, with an apnea-hypopnea index (AHI) $\ge$ 5 events per hour accompanied by symptoms, or AHI $\ge$ 15 regardless of symptoms. Central sleep apnea involves absent respiratory effort due to loss of ventilatory drive, while mixed episodes share features of both.

\section*{What Happens in Your Body}
Obstructive sleep apnea (OSA) occurs when pharyngeal dilator muscles (genioglossus, tensor veli palatini) fail to maintain airway patency during sleep, causing repeated airway collapse at the retropalatal and retroglossal levels. Each apnea triggers progressive low oxygen levels, hypercapnia, and increasing respiratory effort until arousal restores tone and airflow. Cycles of desaturation-reoxygenation activate sympathetic nervous system, increase oxidative stress, and promote systemic inflammation leading to cardiovascular complications.

\section*{Causes}
\begin{itemize}
  \item Obstructive sleep apnea (OSA, most common): Obesity (BMI $>$ 30), craniofacial abnormalities (retrognathia, macroglossia), tonsillar hypertrophy, nasal blockage, hypothyroidism, acromegaly, neuromuscular disorders, menopause, supine sleep position
  \item Central sleep apnea: Heart failure (Cheyne-Stokes respiration), stroke, opioid use, high-altitude periodic breathing, unknown cause central sleep apnea, complex sleep apnea syndrome (treatment-emergent)
  \item Contributing factors: Alcohol use before bedtime, sedative-hypnotic medications, smoking, family history of OSA, male sex and postmenopausal female
\end{itemize}

\section*{When to See a Doctor}
See your doctor if:
\begin{itemize}
  \item Loud, disruptive snoring with witnessed apneas or gasping/choking episodes during sleep
  \item Excessive daytime sleepiness (Epworth Sleepiness Scale score $\ge$ 10) despite adequate sleep time
  \item Morning headaches, dry mouth, or sore throat suggestive of nocturnal mouth breathing
  \item waking at night to urinate (2+ voids/night) without other urologic explanation
  \item Unexplained hypertension (especially resistant), atrial fibrillation, or heart failure
  \item Motor vehicle accident or near-accident attributed to drowsy driving
\end{itemize}

\section*{Related Symptoms}
\begin{itemize}
  \item Snoring
  \item Witnessed apneas
  \item Gasping or choking at night
  \item Excessive daytime sleepiness
  \item Morning headache
  \item waking at night to urinate
  \item Insomnia (sleep-maintenance)
  \item Fatigue
  \item Cognitive impairment
  \item Mood disturbance (depression, irritability)
\end{itemize}
\textbf{Important:} Polysomnography (Level I in-lab or Level III home sleep apnea test) is required for diagnosis; AHI categorizes severity as mild (5--14), moderate (15--30), or severe ($>$ 30 events/hour).

\section*{Self-Care / Home Management}
\begin{itemize}
  \item Weight loss of 10\% or more can significantly reduce AHI and may resolve mild OSA.
  \item Avoid alcohol and sedatives within 3 hours of bedtime.
  \item Sleep on the side rather than supine using positional therapy devices or a tennis ball sewn into pajama back.
  \item Maintain regular sleep schedule with adequate duration (7--9 hours).
  \item Treat nasal congestion with saline sprays or intranasal corticosteroids.
\end{itemize}

\section*{How Your Doctor Will Evaluate You}
\begin{itemize}
  \item History includes snoring volume, witnessed apneas, daytime sleepiness (Epworth scale), sleep partner reports, and Epworth Sleepiness Scale.
  \item Physical examination includes BMI, neck circumference ($>$ 17 inches men, $>$ 16 inches women), Mallampati score, tonsil size, nasal patency, and blood pressure measurement.
  \item Home sleep apnea test (HSAT) is appropriate for high-risk people without comorbidities; polysomnography (PSG) is indicated for those with cardiopulmonary disease or suspected central/complex apnea.
  \item Additional testing includes TSH to rule out hypothyroidism and ECG for arrhythmia screening.
\end{itemize}

\section*{Treatment Options}
\begin{itemize}
  \item Positive airway pressure (PAP) therapy (CPAP, APAP, or BiPAP) is first-line treatment for moderate to severe OSA (AHI $\ge$ 15).
  \item Oral appliance therapy (mandibular advancement device) is an alternative for mild to moderate OSA or people who cannot tolerate PAP.
  \item Lifestyle modification (weight loss, positional therapy, avoidance of alcohol) is recommended for all severity levels.
  \item Upper airway surgery (uvulopalatopharyngoplasty, hypoglossal nerve stimulation, tonsillectomy, or maxillomandibular advancement) is reserved for people with anatomical blockage who fail PAP.
  \item Central sleep apnea is managed by treating the underlying condition (optimize heart failure, discontinue opioids) and may require adaptive servo-ventilation or supplemental oxygen.
\end{itemize}

\clearpage
